% !TEX root = ../algebra_02.tex

\section{Конечность расширения и конечная порождённость}

\begin{Theorem}{Критерий конечности расширения}{}
	Для расширения $L/K$ следующие условия эквивалентны:
	\begin{enumerate}
		\item $L/K$ конечно;
		\item $L/K$ конечно порождено и алгебраично.
	\end{enumerate}
\end{Theorem}

\begin{proof}
	\textbf{$(1 \implies 2)$:} Пусть расширение $L/K$ конечно. Тогда оно имеет конечный базис $a_1, \dots, a_n$ как векторное пространство над $K$. Рассмотрим поле $K_1 = K(a_1, \dots, a_n)$. Так как $K_1$ содержит все $a_i$ и поле $K$, оно обязано содержать все их линейные комбинации:
	\[ \forall \alpha_1, \dots, \alpha_n \in K: \alpha_1 a_1 + \dots + \alpha_n a_n \in K_1 \]
	Но эти линейные комбинации покрывают всё пространство $L$, следовательно, $L = K_1 = K(a_1, \dots, a_n)$. Таким образом, расширение конечно порождено. Алгебраичность любого конечного расширения следует из того, что для любого элемента $a \in L$ степени $1, a, a^2, \dots, a^{[L:K]}$ линейно зависимы, что даёт корень многочлена.

	\textbf{$(2 \implies 1)$:} Пусть $L = K(a_1, \dots, a_n)$, где все $a_i$ алгебраичны над $K$. Докажем конечность расширения индукцией по количеству порождающих элементов $j$. Обозначим $L_j := K(a_1, \dots, a_j)$ для $j = 0, 1, \dots, n$.
	\begin{itemize}
		\item \textit{База:} При $j=0$ имеем $L_0 = K$, и степень расширения $[K:K] = 1 < \infty$.
		\item \textit{Шаг:} Пусть для $j-1$ доказано, что $[L_{j-1}:K] < \infty$. Рассмотрим поле $L_j = L_{j-1}(a_j)$. Элемент $a_j$ алгебраичен над $K$, а значит, подавно алгебраичен и над более широким полем $L_{j-1}$. Поэтому степень простого расширения $[L_j:L_{j-1}] \le \deg_K(a_j) < \infty$. По теореме о башне конечных расширений:
		\[ [L_j:K] = [L_j:L_{j-1}] \cdot [L_{j-1}:K] < \infty \]
	\end{itemize}
	Применяя индукционный переход вплоть до $j=n$, получаем $[L_n:K] = [L:K] < \infty$.
\end{proof}

\begin{Proposition}{Поле алгебраических чисел}{}
	Пусть дано произвольное расширение $L/K$. Множество всех элементов поля $L$, алгебраических над $K$:
	\[ M = \{ x \in L \mid x \text{ алгебраичен над } K \} \]
	является подполем в $L$.
\end{Proposition}

\begin{proof}
	Необходимо проверить, что множество $M$ содержит $0$ и $1$, а также замкнуто относительно сложения, умножения и взятия обратного элемента.
	Очевидно, что $0, 1 \in M$, так как элементы базового поля $K$ являются корнями многочленов первой степени (алгебраичны).

	Пусть $x, y \in M$, причем $x, y \neq 0$. Рассмотрим промежуточное поле $K_1 = K(x, y)$. Так как $x$ и $y$ алгебраичны над $K$, расширение $K_1/K$ является конечно порожденным алгебраическим.
	По критерию конечности, $[K_1:K] < \infty$.

	Любое конечное расширение является алгебраическим, поэтому все элементы поля $K_1$ алгебраичны над $K$. Так как элементы $x+y$, $x \cdot y$ и $x^{-1}$ принадлежат $K_1$, они также алгебраичны над $K$ и, следовательно, лежат в $M$.
	Это доказывает, что $M$ --- поле.
\end{proof}
